% TOP_PROBLEM: {"id": 8400003, "title": "O4 — Finding a prime in an arithmetic progression", "category_id": 3, "category": {"id": 3, "name": "graph_theory", "display_name": "Graph Theory", "description": "Problems involving graphs, networks, and their properties.", "slug": "graph-theory", "order_index": 3, "created_at": "2026-07-31T15:26:25.671Z"}, "status": "open", "problem_number": "AMR-083-0003", "set_id": 13}
% FIRST_POSTED: 2026-10-01
% ENTRY_KIND: statement_only
% UnsolvedMath ID 8400003; source catalogue code AMR-083-0003
% !TeX program = lualatex
% Definition and sources only; this file is not an LLM solution attempt.
\documentclass[11pt,a4paper]{article}
\usepackage[margin=25mm]{geometry}
\usepackage{fontspec}
\setmainfont{lmroman10-regular.otf}[BoldFont=lmroman10-bold.otf,
 ItalicFont=lmroman10-italic.otf,BoldItalicFont=lmroman10-bolditalic.otf]
\usepackage{amsmath,amssymb,mathtools}
\usepackage{unicode-math}
\setmathfont{latinmodern-math.otf}
\AtBeginDocument{\let\setminus\smallsetminus}
\usepackage{xurl,hyperref}
\hypersetup{colorlinks=true,urlcolor=blue}
\setcounter{secnumdepth}{0}
\setlength{\parindent}{0pt}
\setlength{\parskip}{5pt}
\setlength{\emergencystretch}{3em}
\begin{document}
\section*{O4 — Finding a prime in an arithmetic progression}\label{8400003}
{\small UnsolvedMath ID 8400003 · AMR-083-0003 · Graph Theory · AMR Open Problem Lists}
\subsection{Definitions and mathematical statement}
Let $a,n$ be positive integers with $\gcd(a,n)=1$, supplied in binary. An integer $p$ lies in the prescribed arithmetic progression if $p\equiv a\pmod n$, equivalently if $n$ divides $p-a$. We seek a prime $p$ in this residue class; primes dividing $n$ cannot qualify when $n>1$.

Is there a deterministic classical algorithm which, on every such pair $(a,n)$, outputs a prime $p\equiv a\pmod n$ in a number of bit operations bounded by a fixed polynomial in $\ell(a)+\ell(n)$? Here $\ell(t)=\lfloor\log_2t\rfloor+1$, and output writing is charged to the running time. One may first reduce $a$ modulo $n$; for $n=1$ any prime qualifies. There is no requirement that the output be the least prime in the progression.

The question asks for a single uniform algorithm, with no randomness, auxiliary factorization, or unproved analytic assumption. The coprimality hypothesis supplies the natural domain on which a prime is guaranteed to exist. Existence of primes in the progression alone does not provide the required efficient construction: a search through every candidate up to a bound polynomial in $n$ may still take exponential time in the binary input length. Both locating a candidate and establishing that the reported integer is prime belong to the algorithmic task.
\subsection{Short English statement}
Can a prime in any prescribed coprime residue class be constructed deterministically in time polynomial in the binary input size?
\subsection{Sources}
\begin{itemize}
\item[S1] ULAM.AI and the UnsolvedMath Contributors, supplied problems.json, record 8400003 (AMR-083-0003), AMR Open Problem Lists. Catalogue identity and source transcription; CC BY 4.0.
\url{https://huggingface.co/datasets/ulamai/UnsolvedMath}.
\item[S2] Adleman \& McCurley - Open Problems in Number Theory Complexity II (1994), O4, p6 / LNCS p296. Original source indexed by AMR-083-0003.
\url{https://mccurley.org/papers/open.ps.gz}.
\end{itemize}
\end{document}
